% ch-a6-04.tex —— 第 6 篇（Feng–Harris）续：PDF 520–531 / 印刷 507–518
% 本片自印刷 p.507（PDF 520）页首起译；上一片 ch-a6-03.tex 覆盖印刷 495–506。
% 页末截断句按原书页界断开；下一片 ch-a6-05.tex 从印刷 p.519（PDF 532）续。

\ysrc{507}

上文提到的诸猜想关乎附着于自守形式的伽罗瓦表示的存在性，以及这类表示的局部–整体相容性。这些假设是剩余表示上通常「Taylor–Wiles」型条件的一种严格形式。\footnote{从技术角度看，[GV18] 应有很大的优化余地，办法是纳入 Kisin 的方法。}确切表述见 [GV18, §6.6, §10]。

\subsubsection{与 Taylor–Wiles 方法的联系}

Galatius–Venkatesh 主结果的证明基于 Taylor–Wiles 方法，并纳入了 Calegari–Geraghty 的修正。一个有趣之处是：$\pi_*(\mathcal R_0)$ 的确定似乎与 Taylor–Wiles 方法紧密相连。下面给一个简短、且极其印象式的勾勒。

先作一点心理定位。我们认为应当存在 $\mathcal R_0$ 对 $\mathcal M_0$ 的一个作用；但在此处及全篇，我们只能构造这一作用的「经典影子」（classical shadow），至少一开始如此，如下图所示：
\[
\begin{tikzcd}[column sep=4.5em, row sep=3em]
\mathcal R_0 \arrow[r, bend left=28, "?"] \arrow[d, Rightarrow] & \mathcal M_0 \arrow[d, Rightarrow] \\
\mathrm R_0 \arrow[r, bend left=28] & M_0
\end{tikzcd}
\]
这里带 ? 的 $\curvearrowright$ 表示一种我们认为存在、但尚不知如何构造的作用。

Taylor–Wiles 方法的出发点是考虑辅助的一族类似情形：在精心选取的一组素数 $Q_n$ 处添加额外的水平结构。设 $Q_n$ 为这样一组素数，我们作如下定义。
\begin{itemize}[leftmargin=2em]
\item $\mathcal R_n := \mathcal R^{\bar\rho,\mathrm{crys}}_{\mathcal O_F[1/SQ_n]}$，$\mathrm R_n := \pi_0\bigl(\mathcal R^{\bar\rho,\mathrm{crys}}_{\mathcal O_F[1/SQ_n]}\bigr)$；由 §5.3.2，后者是经典晶体形变环。导出环 $\mathcal R_n$ 的虚拟维数为 $-\delta$。
\item $\mathcal M_n$ 为对应于「$G$ 的水平 $SQ_n$ 的自守形式」的模；按假设，它是集中在次数 $[-\delta,0]$ 的复形。令 $M_n := H_0(\mathcal M_n)$。
\end{itemize}

于是，对每个 $\mathcal R_n$ 都有同样的故事：我们想要一个作用 $\mathcal R_n\overset{?}{\curvearrowright}\mathcal M_n$，但实际能构造的只是它的经典影子 $\mathrm R_n\curvearrowright M_n$。因此有图
\[
\begin{tikzcd}[column sep=4.5em, row sep=3em]
\mathcal R_n \arrow[r, bend left=28, "?"] \arrow[d, Rightarrow] & \mathcal M_n \arrow[d, Rightarrow] \\
\mathrm R_n \arrow[r, bend left=28] & M_n
\end{tikzcd}
\]

初看之下，图景与 $n=0$ 时毫无二致，似乎没有进展。但要点在于：巧妙选取 $Q_n$，可以在道德上让 $\mathcal R_n$「趋于」一个实际上是经典的对象 $\mathcal R_\infty$，使得 $\mathcal R_\infty\xrightarrow{\sim}\mathrm R_\infty:=\pi_0(\mathcal R_\infty)$。现实中我们不会这样构造 $\mathcal R_\infty$；我们要构造的是一个「光滑且维数正确」的 $\mathrm R_\infty$。由 §3.2.1 的讨论，这为令 $\mathcal R_\infty$ 等于 $\mathrm R_\infty$ 提供了理由。

进而，$\mathcal M_n$「趋于」一个对象 $\mathcal M_\infty$（它确实存在），并且确有 $\mathcal M_\infty\xrightarrow{\sim}M_\infty$。还应设想：图 (5.5.3)「取极限」后得到下图：\tn{底本此处称「图 (5.5.3)」，但 (5.5.3) 给出的是公式而非图，疑为作者笔误（原书 p.507–508）。}

\ysrc{508}

\[
\begin{tikzcd}[column sep=4.5em, row sep=3em]
\mathcal R_\infty \arrow[r, bend left=28, "?"] \arrow[d, Rightarrow] & \mathcal M_\infty \arrow[d, Rightarrow] \\
\mathrm R_\infty \arrow[r, bend left=28] & M_\infty
\end{tikzcd}
\]
此情形下，$\mathcal R_\infty\onto\mathrm R_\infty$ 实为单纯交换环的同构，故它自动给出 $\mathcal M_\infty\xrightarrow{\sim}M_\infty$ 上所想要的作用。

为看清这如何有助于原问题，我们来考察 $\mathcal R_n$ 与 $\mathcal R_0$ 的关系。差别在于：$\mathcal R_0$ 只参数化在 $Q_n$ 处非分歧的形变（回忆 $\bar\rho$ 在 $Q_n$ 中各素处非分歧），而 $\mathcal R_n$ 允许在 $Q_n$ 处分歧的形变。如例 5.4.1 所述，可以从 $\mathcal R_n$ 出发作一个导出张量积，从 $Q_n$ 上各局部伽罗瓦群的形变中截出非分歧轨迹，从而复原 $\mathcal R_0$。更精确地说，选取 $Q_n$ 使得 $\bar\rho$ 的形变在 $Q_n$ 上方自动驯顺。于是，从 $Q_n$ 上的驯顺惯性群到 $\pi_1(\mathcal O_F[1/SQ_n])$ 的映射诱导出环同态 $S_n\to\mathcal R_n$，其中 $S_n$ 是 $Q_n$ 中各素处的驯顺惯性的某个形变环。存在增广 $S_n\to S_0=W(k)$，对应那些在驯顺惯性上平凡的形变。与例 5.4.1 中类似的条件给出单纯交换环的等价
\begin{equation}\tag{5.5.1}\label{eq:a6-04-1}
\mathcal R_n\overset{\mathbf L}{\otimes}_{S_n}S_0\xrightarrow{\sim}\mathcal R_0.
\end{equation}

此恒等式近似地「取极限」为
\begin{equation}\tag{5.5.2}\label{eq:a6-04-2}
\text{「}\mathcal R_\infty\overset{\mathbf L}{\otimes}_{S_\infty}S_0\xrightarrow{\sim}\mathcal R_0.\text{」}
\end{equation}

更一般地，近似地有
\begin{equation}\tag{5.5.3}\label{eq:a6-04-3}
\text{「}\mathcal R_\infty\overset{\mathbf L}{\otimes}_{S_\infty}S_n\xrightarrow{\sim}\mathcal R_n.\text{」}
\end{equation}

我们之所以给这些陈述加上引号，是因为它们并不精确成立，而是差一个误差；该误差随 $n\to\infty$ 趋于 $0$，因此对我们的目的而言可以接受。这一「误差」的来源在于 $\mathrm R_\infty$ 的构造：它由 $\mathrm R_n$ 的有限商的逆向系统经紧性论证抽出。

此外还有 $S_n$ 通过「菱形算子」（diamond operators）对 $M_n$ 的作用，对此有
\begin{equation}\tag{5.5.4}\label{eq:a6-04-4}
\mathcal M_\infty\overset{\mathbf L}{\otimes}_{S_\infty}S_0\xrightarrow{\sim}\mathcal M_0.
\end{equation}

再者，局部–整体相容性保证这一作用与 $S_n\to\mathrm R_n$ 所诱导的作用相容。于是，由 $\mathcal R_\infty\xrightarrow{\sim}\mathrm R_\infty$ 对 $M_\infty\xleftarrow{\sim}\mathcal M_\infty$ 的作用，取导出张量积并用 (5.5.2) 与 (5.5.4)，得到
\[
\text{「}\mathcal R_0\cong\mathcal R_\infty\overset{\mathbf L}{\otimes}_{S_\infty}S_0\text{」}
\quad\curvearrowright\quad
\text{「}\mathcal M_\infty\overset{\mathbf L}{\otimes}_{S_\infty}S_0\cong\mathcal M_0\text{」}.
\]

如上所述，加引号的陈述并不精确，但可以在同伦群上做到任意精度。办法是改用 (5.5.3)「从 $\infty$ 走到 $n$」，其引入的误差随 $n$ 增大而任意小；再用 (5.5.1)「从 $n$ 走到 $0$」，则不引入任何误差。这足以给出一个同构

\ysrc{509}

\[
\pi_i(\mathcal R_0)\approx\pi_i\bigl(\mathcal R_\infty\overset{\mathbf L}{\otimes}_{S_\infty}S_0\bigr)\cong\mathrm{Tor}^{S_\infty}_i(\mathcal R_\infty,S_0),
\]
随后我们便能（在有利的假设下）显式地计算它。事实上，上述论证还表明：$S_\infty\to\mathrm R_\infty$ 是模一个长为 $\delta$ 的正则序列的商；由此可用标准的 Koszul 复形计算 $\mathrm{Tor}^{S_\infty}_i(\mathcal R_\infty,S_0)$。

至此，这段毫不掩饰其含糊与印象式的勾勒就结束了。建议读者参阅 Caraiani–Shin [CS25] 的文章，其中对 Calegari–Geraghty 修正版的 Taylor–Wiles 方法有实质得多的讨论。

\section{导出 Hecke 代数}

\subsection{局部导出 Hecke 代数}

我们简要回顾 [Ven19] 中的导出 Hecke 代数理论。设 $G$ 是局部域 $K$ 上的分裂约化群，其剩余域为 $\F_q$，特征为 $p$。设 $U\subset G(K)$ 为紧开子群。（就我们的目的而言，最关心的是极大紧子群 $U=G(\mathcal O_K)$。）

设 $\Lambda$ 为交换环。定义泛模
\[
\mathfrak U_\Lambda(U)=\mathrm{c\text{-}Ind}^{G(K)}_U\Lambda,
\]
其中 $U$ 在 $\Lambda$ 上平凡作用。$(G(K),U)$ 的通常 Hecke 代数可以表为
\[
\mathrm H(G(K),U;\Lambda):=\mathrm{Hom}_{G(K)}\bigl(\mathfrak U_\Lambda(U),\mathfrak U_\Lambda(U)\bigr).
\]
这一表述提示如下推广。

\begin{defn}{6.1.1}\label{thm:a6-04-1}
$(G(K),U)$ 的系数在 $\Lambda$ 中的\emph{导出 Hecke 代数}是分次环
\[
\mathcal H(G(K),U;\Lambda):=\mathrm{Ext}^*_{G(K)}\bigl(\mathfrak U_\Lambda(U),\mathfrak U_\Lambda(U)\bigr),
\]
其中 $\mathrm{Ext}$ 在光滑 $G(K)$-表示范畴中取。当 $U=G(\mathcal O_K)$ 时，简记 $\mathcal H(G(K);\Lambda):=\mathcal H(G(K),G(\mathcal O_K);\Lambda)$。
\end{defn}

\begin{rem}{6.1.2}\label{thm:a6-04-2}
环 $\mathcal H(G(K),U;\Lambda)$ 其实只是一个分次（结合）环，而非任何意义下的导出（结合）环。不过，它的描述清楚地表明：它是微分分次代数 $\mathrm{RHom}_{G(K)}(\mathfrak U_\Lambda(U),\mathfrak U_\Lambda(U))$ 的上同调群。更清教徒式的做法或许是把后者称为「导出 Hecke 代数」，但我们沿用 [Ven19]，把这个名字用于分次环 $\mathcal H(G(K),U;\Lambda)$。我们将改称 $\mathrm{RHom}_{G(K)}(\mathfrak U_\Lambda(U),\mathfrak U_\Lambda(U))$ 为\emph{微分分次 Hecke 代数}，记作 $\mathscr H(G(K),U;\Lambda)$；当 $U=G(\mathcal O_K)$ 时有时简记为 $\mathscr H(G(K);\Lambda)$。它没有动态环（animated ring）的结构；先验地，它只是一个微分分次（即 $\mathrm E_1$）代数，尽管我们在 §9.2 猜想它具有更多的交换结构。
\end{rem}

\begin{exam}{6.1.3}\label{thm:a6-04-3}
$\mathcal H(G(K),U;\Lambda)$ 的第 $0$ 个分次群是经典 Hecke 代数 $\mathrm H(G(K),U;\Lambda)$。
\end{exam}

下面按 [Ven19, §2] 给出导出 Hecke 代数的另外两个更具体的描述。

\ysrc{510}

\subsubsection{函数论描述}

设 $x,y\in G(K)/U$，$G_{x,y}\subset G(K)$ 为配对 $(x,y)$ 的稳定化子。注意 $G_{x,y}$ 是 $G(K)$ 的紧开子群。我们可以把 $\mathcal H(G(K),U;\Lambda)$ 看作由满足下列条件的函数构成的空间：
\[
G(K)/U\times G(K)/U\ni(x,y)\mapsto h(x,y)\in H^*(G_{x,y};\Lambda).
\]
\begin{enumerate}[(1)]
\item 函数 $h$ 在左边「$G(K)$-不变」。更精确地说，对一切 $g\in G(K)$ 有
\[
[g]^*h(gx,gy)=h(x,y),
\]
其中 $[g]^*\colon H^*(G_{gx,gy};\Lambda)\to H^*(G_{x,y};\Lambda)$ 是沿 $\mathrm{Ad}(g)$ 的拉回。
\item $h$ 的支撑模 $G(K)$ 有限。
\end{enumerate}
乘法由卷积公式给出：在余定义域上用杯积定义乘法，并用限制/膨胀把上同调类移到正确的群（见 [Ven19, eqn. (22)]）。

\subsubsection{双陪集描述}

对 $x\in G(K)/U$，令 $U_x=\mathrm{Stab}_U(x)$。具体地说，若 $x=g_xU$，则 $U_x:=U\cap g_xUg_x^{-1}$。

我们也可以把 $\mathcal H(G(K),U;\Lambda)$ 描述为紧支撑（即支撑于有限多个双陪集上）的函数
\[
x\in U\backslash G(K)/U\mapsto h(x)\in H^*(U_x;\Lambda).
\]
至于代数结构的描述，改用 6.1.1 的模型似乎更好。

\subsubsection{环面的导出 Hecke 代数}

设 $T$ 为 $K$ 上的分裂环面。现在利用双陪集模型，可以显式地描述环面 $T(K)$ 的导出 Hecke 代数。由于 $T$ 交换，对一切 $x\in T(K)/T(\mathcal O_K)$ 都有 $T(\mathcal O_K)_x=T(\mathcal O_K)$。

我们有 $T(K)/T(\mathcal O_K)\cong X_*(T)$。作等同
\begin{equation}\tag{6.1.1}\label{eq:a6-04-5}
X_*(T)=T(K)/T(\mathcal O_K)\hookrightarrow G(K)/G(\mathcal O_K)
\end{equation}
所用的映射为 $X_*(T)\ni\chi\mapsto\chi(\varpi)\in G(K)/G(\mathcal O_K)$，其中 $\varpi$ 是 $K$ 的一致化子。于是 $\mathcal H(T(K);\Lambda)$ 恰好由紧支撑函数
\[
X_*(T)\to H^*(T(\mathcal O_K);\Lambda)
\]
组成，乘法由卷积给出；换言之，
\[
\mathcal H(T(K);\Lambda)\cong\Lambda[X_*(T)]\otimes_\Lambda H^*(T(\mathcal O_K);\Lambda).
\]

\begin{rem}{6.1.4}\label{thm:a6-04-4}
若 $p$ 在 $\Lambda$ 中可逆，则 $H^*(T(\mathcal O_K);\Lambda)$ 的结构可如下阐明。首先，约化映射 $T(\mathcal O_K)\to T(\F_q)$ 的核是 pro-$p$ 群；按假设，它以 $\Lambda$ 为系数的高阶上同调消没，故诱导出同构 $H^*(T(\F_q);\Lambda)\to H^*(T(\mathcal O_K);\Lambda)$。

其次，因 $T$ 分裂，有 $T(\F_q)\cong(\F_q^\times)^r$，其中 $r=\mathrm{rank}(T)$。我们主要关心 $\Lambda$ 形如 $\mathcal O/l^n$ 的情形，其中 $\mathcal O$ 是数域的整数环，$l$ 是 $\ell\ne p$ 之上的素元。此时，$H^*(T(\F_q);\Lambda)$ 的高阶上同调不消没，当且仅当 $q\equiv1\pmod\ell$。
\end{rem}

\ysrc{511}

\subsubsection{导出 Satake 同构}

本小节假设 $q\equiv1\in\Lambda$。设 $T\subset G$ 为极大环面，$W$ 为 $T\subset G$ 的 Weyl 群。我们考虑导出 Hecke 代数 $\mathcal H(G(K);\Lambda)$ 的经典 Satake 变换的一个类比，其形式为
\begin{center}
\itshape ``Derived Hecke algebra for $G\isom$ (Derived Hecke algebra for $T$)$^W$.''
\end{center}

更精确地说，我们定义\emph{导出 Satake 变换}
\begin{equation}\tag{6.1.2}\label{eq:a6-04-6}
\mathcal H(G(K);\Lambda)\to\mathcal H(T(K);\Lambda)
\end{equation}
为沿映射
\[
\bigl(T(K)/T(\mathcal O_K)\bigr)^2\to\bigl(G(K)/G(\mathcal O_K)\bigr)^2
\]
的限制（用 §6.1.1 的函数论模型），该映射来自 (6.1.1)。更详细地说，设 $h\in\mathcal H(G(K);\Lambda)$ 由函数
\[
\bigl(G(K)/G(\mathcal O_K)\bigr)^2\ni(x,y)\mapsto h(x,y)\in H^*(G_{x,y};\Lambda)
\]
给出。则 (6.1.2) 把 $h$ 送到复合
\[
\bigl(T(K)/T(\mathcal O_K)\bigr)^2\longleftarrow\bigl(G(K)/G(\mathcal O_K)\bigr)^2
\xrightarrow{\ h\ }H^*(G_{x,y};\Lambda)\xrightarrow{\ \mathrm{res}\ }H^*(T_{x,y};\Lambda).
\]

\begin{rem}{6.1.5}\label{thm:a6-04-5}
这个定义是正确的，这一点或许令人惊讶：因为在特征 $0$ 下，通常（非导出）Hecke 代数上的类比构造远非通常的 Satake 变换。只是由于我们对特征关系的假设（即 $q\equiv1\in\Lambda$），这一「朴素」（naïve）构造才表现良好；否则它连一个环同态都不是。该构造的动机似乎来自 [TV16]；相关要点的证明见 [Fen24, Lemma 6.5]。
\end{rem}

\begin{thm}{6.1.6}（[Ven19, Theorem 3.3]）\label{thm:a6-04-6}
设 $W$ 为 $T$ 在 $G$ 中的 Weyl 群。若 $|W|$ 在 $\Lambda$ 中可逆且 $q=1\in\Lambda$，则映射 (6.1.2) 诱导同构
\[
\mathcal H(G(K);\Lambda)\xrightarrow{\sim}\mathcal H(T(K);\Lambda)^W.
\]
\end{thm}

\begin{rem}{6.1.7}\label{thm:a6-04-7}
严格说来，[Ven19] 只在 $K$ 的特征为零时证明了定理 6.1.6，但本质上同样的论证应当普遍成立。
\end{rem}

\subsubsection{交换性？}

先验地，$\mathscr H(G(K),U;\Lambda)$ 只是一个结合微分分次代数，$\mathcal H(G(K),U;\Lambda)$ 只是一个结合分次代数。有趣的问题是：这两个代数还具有哪些更多的结构或性质。事实上，$H^0(\mathscr H(G(K),U;\Lambda))\cong\mathcal H(G(K),U;\Lambda)$ 众所周知是交换的，尽管要证明这一点并不平凡。

我们认为，猜想 $\mathcal H(G(K);\Lambda)$ 也交换是合理的；至少当 $\Lambda$ 的特征相对于 $G$ 与 $K$ 的剩余特征不处于「坏」位置时如此。（目前即使对坏特征也没有反例。）我们列出若干已知的交换性情形：
\begin{enumerate}[(1)]
\item 当 $\Lambda$ 的特征 $\ell\nmid\#W$ 且 $q\equiv1\pmod\ell$ 时，由定理 6.1.6 知交换性成立（该结果属于 Venkatesh）。
\item 更一般地，若 $\#G/B(\F_q)$ 与 $p$ 都在 $\Lambda$ 中可逆，则由 Gehrmann [Geh22] 知交换性成立。\footnote{该文还考察了 $G=\mathrm{SL}_2$ 的一些不满足这些假设的例子。}
\end{enumerate}

\ysrc{512}

为对这些论证略作说明，我们回忆经典球 Hecke 代数交换性的两个证明：一个经由 Satake 变换，一个经由「Gelfand 技巧」。Venkatesh 的论证是 Satake 变换方法的推广，Gehrmann 的论证是 Gelfand 技巧的推广。我们注意到：虽然 Gehrmann 证明的交换性更为一般，但 Venkatesh 的证明给出了整体应用所需的更精细信息。

除 $\mathcal H(G(K);\Lambda)$ 的交换性之外，还可以提出更微妙的问题，它们涉及微分分次 Hecke 代数 $\mathscr H(G(K);\Lambda)$ 的结构。§9.2 将对此问题作一些更具推测性的讨论。

\subsubsection{推广}

关于导出 Hecke 代数的文献尚不如人们所期望的那样全面。以下推广会是有用的：
\begin{itemize}[leftmargin=2em]
\item 把理论推广到非分裂群 $G$。
\item 把理论推广以纳入非平凡系数，即把 $\mathfrak U_\Lambda(U)$ 换成 $U$ 的某个非平凡表示的紧诱导。在整体情形下，这与在相应的局部对称空间上考虑非平凡局部系统数有关。
\end{itemize}

\subsection{导出 Hecke 代数的作用}

沿用前面的记号：设 $G$ 是剩余特征为 $p$ 的局部域 $K$ 上的约化群，$U\subset G(K)$ 为紧开子群，$\Lambda$ 为交换环。

\begin{fact}{6.2.1}\label{thm:a6-04-8}
若 $C^\bullet$ 是 $\Lambda$ 上光滑 $G(K)$-表示的导出范畴的一个对象，则（形式上）微分分次 Hecke 代数 $\mathscr H(G(K),U;\Lambda)$（相应地，导出 Hecke 代数 $\mathcal H(G(K),U;\Lambda)$）右作用在 $\mathrm{RHom}_{\Lambda[U]}(\Lambda,C^\bullet)$（相应地，$\mathrm{Ext}^*_{\Lambda[U]}(\Lambda,C^\bullet)$）上。
\end{fact}

事实上，导出自同态环 $\mathrm{RHom}_{G(K)}(\mathfrak U_\Lambda(U),\mathfrak U_\Lambda(U))$ 重言地右作用在函子 $\mathrm{RHom}^\bullet_{\Lambda[G(K)]}(\mathfrak U_\Lambda(U),-)$ 上。于是该事实由 Frobenius 互反律同构
\begin{equation}\tag{6.2.1}\label{eq:a6-04-7}
\mathrm{RHom}^\bullet_{\Lambda[U]}(\Lambda,C^\bullet)\xrightarrow{\sim}\mathrm{RHom}^\bullet_{\Lambda[G(K)]}\bigl(\mathfrak U_\Lambda(U),C^\bullet\bigr)
\end{equation}
得出。

函子 $C^\bullet\mapsto\mathrm{RHom}^\bullet_{\Lambda[U]}(\Lambda,C^\bullet)$ 就是导出 $U$-不变量。因此，事实 6.2.1 可用如下标语更通俗地概括：
\begin{center}
\itshape There is a canonical right action of the derived Hecke algebra for $(G,U)$ on the derived $U$-invariants of a $G(K)$-representation.
\end{center}

\begin{rem}{6.2.2}\label{thm:a6-04-9}
若 $U$ 是 pro-$p$ 群且 $p$ 在 $\Lambda$ 中可逆，则取 $U$-不变量的操作正合。$p$ 在 $\Lambda$ 中不可逆的情形将在下一小节讨论。
\end{rem}

\begin{exam}{6.2.3}\label{thm:a6-04-10}
例子的一类来源是局部对称空间的上同调。取 $U$ 为 $p$ 处的水平结构，则 $G$ 的局部对称空间的上同调复形可典范地实现为某个 $G(K)$-表示的导出 $U$-不变量；这一表示至少启发式地是「$p$ 处具有无限水平结构的局部对称空间的上同调复形」。[Ven19] 中局部导出 Hecke 代数的动机正在于分析此类例子。后文各节还会更多地讨论这一点。
\end{exam}

\ysrc{513}

\subsection{$p$-进群的 mod $p$ 导出 Hecke 代数}

设 $G$ 是 $\Q_p$ 上的约化群——我们对代数群与其 $\Q_p$-点群使用同一记号——并设 $k$ 是特征为 $p$ 的代数闭域。记 $\mathrm{Rep}_k(G)$ 为 $G$ 在 $k$-向量空间上连续表示的范畴，$\mathrm{Rep}^{\mathrm{sm}}_k(G)\subset\mathrm{Rep}_k(G)$ 为光滑表示的子范畴。当 $G=\mathrm{GL}_2(\Q_p)$ 或与之密切相关的群时，$\mathrm{Rep}^{\mathrm{sm}}_k(G)$ 中的不可约对象早已得到分类；但对其他任何群，即使只是近似的分类也仍然难以企及，尽管过去 15 年里有 [Bre03, Koz16, BP12] 等令人印象深刻的努力。

尽管如此，阿贝尔范畴 $\mathrm{Rep}^{\mathrm{sm}}_k(G)$ 在一个方面结构简单：每个不可约对象都典范地带有从一个紧投射生成元出发的满同态空间，此生成元下文记作 $\mathfrak U$。Schneider 利用这一事实，在 $\mathrm{Rep}^{\mathrm{sm}}_k(G)$ 与导出自同态代数 $R\mathrm{End}_G(\mathfrak U)^{\mathrm{op}}$ 上的 dg-模范畴之间建立了一个导出等价。后者的性质是上述 mod $\ell$ 导出 Hecke 代数的 mod $p$ 类比，目前仍大多未被探索。一个更简单的版本称为导出菱形代数；Khare 与 Ronchetti [KR23a] 证明它在 $p$-进模形式理论中起着重要作用。

\subsubsection{Schneider 理论}

固定 Iwahori 子群 $I\subset G$，并设 $I(1)\subset I$ 为其极大 pro-$p$ 正规子群；于是 $I/I(1)$ 是某个有限域 $\F$ 上环面的 $\F$-点群。容易证明：每个不可约的 $\pi\subset\mathrm{Rep}^{\mathrm{sm}}_k(G)$ 都由其子空间 $\pi^{I(1)}$ 生成。这定义了一个典范函子
\begin{equation}\tag{6.3.1}\label{eq:a6-04-8}
\begin{gathered}
\tau_k\colon\mathrm{Rep}^{\mathrm{sm}}_k(G)\to\mathrm{Mod}\bigl(\mathrm H(G,I(1);k)^{\mathrm{op}}\bigr)\\
\pi\mapsto\pi^{I(1)}
\end{gathered}
\end{equation}
其中 $\mathrm H(G,I(1);k)$ 是 $G$ 相对 $I(1)$ 的 $k$-值 Hecke 代数。这里要提醒一个记号上的不一致：Schneider 的约定与我们（随 Venkatesh）的约定相差一个反代数。也就是说，令 $\mathfrak U:=\mathrm{c\text{-}Ind}^G_{I(1)}k$ 为泛诱导模，则 Schneider 把 $G$ 相对 $I(1)$ 的 Hecke 代数定义为 $\mathrm{End}_G(\mathfrak U)^{\mathrm{op}}$，而非 $\mathrm{End}_G(\mathfrak U)$。因此，把他的结果翻译到我们的约定中时，会多出一个「op」。

Borel 与 Casselman 的一个著名定理 [Bor76] 断言：若把 $k$ 换成 $\C$，则 (6.3.1) 在由 $I(1)$-不动向量生成的表示所构成的子范畴与 $\mathrm{Mod}(\mathrm H(G,I(1);k)^{\mathrm{op}})$ 之间定义了一个范畴等价。当 $G=\mathrm{GL}_2(\Q_p)$ 且 $k$ 的特征为 $p$ 时，这同样成立；但 Ollivier [Oll09] 证明，对其他 $G$，这几乎从不成立——尽管每个不可约 $\pi$ 都由 $\pi^{I(1)}$ 生成。Schneider 通过把 $\mathrm H(G,I(1);k)$ 换成微分分次代数
\begin{equation}\tag{6.3.2}\label{eq:a6-04-9}
\mathscr H(G,I(1);k)=R\mathrm{End}_G(\mathfrak U):=\mathrm{RHom}_G(\mathfrak U,\mathfrak U).
\end{equation}
找到了 Borel–Casselman 定理的合适推广。函子 (6.3.1) 的替代品是
\begin{equation}\tag{6.3.3}\label{eq:a6-04-10}
\begin{gathered}
\tau^\bullet_k\colon D(G)\to D\bigl(\mathscr H(G,I(1);k)^{\mathrm{op}}\bigr)\\
\pi\mapsto\mathrm{RHom}_{I(1)}(1,\pi)=\mathrm{RHom}_G(\mathfrak U,\pi)
\end{gathered}
\end{equation}
这里 $D$ 表示（无界）导出范畴，我们用 $D(G)$ 代替 $D(\mathrm{Rep}^{\mathrm{sm}}_k(G))$，而最后一个等式是 Frobenius 互反律。Schneider 的定理如下。

\ysrc{514}

\begin{thm}{6.3.1}（[Sch15a], Theorem 9）\label{thm:a6-04-11}
设 $I(1)$ 是无挠的 $p$-进群。则 $\mathfrak U$ 是 $D(G)$ 的紧生成元，且函子 (6.3.3) 是三角范畴的等价。
\end{thm}

关于范畴 $\mathrm{Rep}^{\mathrm{sm}}_k(G)$ 的第一个描述见 [AHHV17]。抛物诱导定义为从 $\mathrm{Rep}^{\mathrm{sm}}_k(M)$ 到 $\mathrm{Rep}^{\mathrm{sm}}_k(G)$ 的函子（其中 $M$ 是抛物子群 $P$ 的 Levi 分量）；在特征 $p$ 下，这一诱导未作归一化，其结果依赖于 $P$ 以及 $M$。若 $G$ 的不可约容许表示 $\pi$ 不作为某个真抛物子群的容许不可约表示所诱导的表示的次商出现，则称之为\emph{超尖点的}。[AHHV17] 的主定理把 $\mathrm{Rep}^{\mathrm{sm}}_k(G)$ 的分类化归为 Levi 子群的超尖点表示的分类。这对应于（非导出的）Hecke 代数 $\mathrm H(G,I(1);k)^{\mathrm{op}}$ 的模的类似分类：对应于超尖点 $\pi$ 的 Hecke 代数模称为\emph{超奇异的}，它们在 $\mathrm H(G,I(1);k)^{\mathrm{op}}$ 中有简单的刻画（见 [AHHV17, p. 498]）。

\begin{rem}{6.3.2}\label{thm:a6-04-12}
这提示：$D(\mathscr H(G,I(1);k)^{\mathrm{op}})$ 的结构也可以化归为超奇异模的三角子范畴的研究；但似乎还没有人研究过这个问题。
\end{rem}

当 $I(1)$ 不是无挠群时，在 $\mathfrak U$ 的定义与定理 6.3.1 中可以用一个适当的有限指数无挠子群来替代它；但一般紧开子群的非导出 Hecke 代数理论要不清楚得多。

与 $\ell\ne p$ 的情形一样，我们把上同调环
\[
\mathcal H(G,U;\Lambda):=H^*(\mathscr H(G,U;\Lambda))
\]
称为\emph{导出 Hecke 代数}；先验地，它只是一个分次结合环。N. Ronchetti 的学位论文 [Ron] 研究了导出球 Hecke 代数 $\mathcal H(G,G(\mathcal O);\Lambda)$，并定义了一个 Satake 同态
\[
\mathcal H(G,G(\mathcal O);\Lambda)\to\mathcal H(T,T(\mathcal O);\Lambda)
\xrightarrow{\sim}\Lambda[X_*(T)]\otimes H^*(T(\mathcal O);\Lambda).
\]
此情形下，若 $T(\mathcal O)$ 无 $p$-挠，则 $H^*(T(\mathcal O),\Lambda)$ 是 $H^1$ 上的外代数。这个代数比 Schneider 的 $\mathscr H(G,I(1);k)$ 容易处理得多；然而，平凡表示在 $G(\mathcal O)$ 上的紧诱导并不生成 $D(G)$，所以 Ronchetti 的代数无法完全解释 $G$ 的 mod $p$ 表示论。不过，更简单的 $p$-进导出 Hecke 代数似乎确实在 $p$-进模形式的整体理论中起作用；见下文 §8.4.2。

\section{局部对称空间的上同调}

\subsection{局部对称空间}

本节中，$G$ 是 $\Q$ 上的连通约化群，中心 $Z=Z_G$。我们固定一个子群 $K_\infty\subset G(\R)$，它包含 $Z(\R)$ 以及恒等分支 $G(\R)^0\subset G(\R)$ 的一个极大紧子群，使得 $K_\infty/Z(\R)$ 紧。则 $X=G(\R)/K_\infty$ 是导子群 $G(\R)^{\mathrm{der},0}$ 的对称空间的有限个拷贝之并。设 $K_f\subset G(\mathbb{A}_f)$ 为紧开子群；$G(\R)$ 中相应的\emph{同余子群}是 $G(\mathbb{A}_f)$ 中的交 $G(\Q)\cap K_f$ 所给出的子群 $\Gamma=\Gamma_{K_f}\subset G(\Q)$。设 $K'_f\subset K_f$ 是有限指数的无挠正规子群，商为 $Q$；则
\[
X(\Gamma_{K'_f}):=\Gamma_{K'_f}\backslash X
\]

\ysrc{515}

它是 $C^\infty$ 流形，且
\[
X(\Gamma_{K_f})=\Gamma_{K_f}\backslash X=Q\backslash X(\Gamma_{K'_f})
\]
是附着于 $K_f$ 的局部对称空间。就多数目的而言，有限群 $Q$ 与 $X(\Gamma_{K_f})$ 的奇点在下文中不起作用。

我们最关心的是阿代尔局部对称空间
\[
S_{K_f}(G,X)=G(\Q)\backslash X\times G(\mathbb{A}_f)/K_f
\]
及其极限
\[
S(G,X)=\varprojlim_{K_f}S_{K_f}(G,X),
\]
其中极限取遍所有紧开子群 $K_f\subset G(\mathbb{A}_f)$。固定 $K_f$ 时，该阿代尔局部对称空间可等同于有限个离散算术商的非交并：由约化理论，可写 $G(\mathbb{A}_f)=\coprod_j G(\Q)\alpha_jK_f$（$\alpha_j$ 取遍一个有限集），于是
\[
S_{K_f}(G,X)=\coprod_j X(\Gamma_j),
\]
其中 $\Gamma_j=G(\Q)\cap\alpha_jK_f\alpha_j^{-1}$。

除另有说明外，对任何上同调理论 $H^*$（通常带扭系数），空间 $H^*(S(G,X))$ 都理解为滤过余极限 $\varinjlim_{K_f}H^*(S_{K_f}(G,X))$。

\subsection{de Rham 理论回顾}

设 $\mathcal A(G)$ 表示 $G(\Q)\backslash G(\mathbb{A})$ 上自守形式构成的（复）向量空间，$\mathcal A_0(G)\subset\mathcal A(G)$ 表示尖点形式的子空间。设 $Z=Z_G$ 为 $G$ 的中心，$K_\infty\subset G(\R)$ 为恒等分支的一个闭子群，它包含 $Z(\R)$，且其在伴随群 $G(\R)^{\mathrm{ad}}$ 中的像紧。视情形而定，$K_\infty$ 可能是、也可能不是模中心的极大紧子群。为记账方便，记 $K^c_\infty$ 为 $K_\infty$ 的极大紧子群。记 $\frakg=\mathrm{Lie}(G)$，$\mathfrak{k}=\mathrm{Lie}(K_\infty)$；除另有说明，均取复点上的李代数。于是有 Cartan 分解
\begin{equation}\tag{7.2.1}\label{eq:a6-04-11}
\frakg=\mathfrak{k}\oplus\frakp.
\end{equation}

（可能不连通的）对称空间 $X=G(\R)/K_\infty$ 具有不变 Hermite 结构，当且仅当存在同态
\[
h\colon\C^\times=R_{\C/\R}\mathbb G_{m,\C}\to G_\R,
\]
其像含于 $K_\infty$ 的中心，且满足志村数据（Shimura datum）的公理（参看本卷 Morel 的文章 [Mor25]）。此时，(7.2.1) 可细化为 $\mathfrak{k}$-不变的 Harish-Chandra 分解：
\begin{equation}\tag{7.2.2}\label{eq:a6-04-12}
\frakg=\mathfrak{k}\oplus\frakp^+\oplus\frakp^-.
\end{equation}
这里，在切空间映射
\[
\frakg=T_{G,e}\to T_{X,h}
\]
下（其中 $T_{G,e}$ 是恒等元处的切空间，$T_{X,h}$ 是 $K_\infty$-不动点 $h\in X$ 处的切空间），$\frakp^+$（相应地 $\frakp^-$）被映到全纯切空间

\ysrc{516}

而 $\mathfrak p^-$ 被映到反全纯切空间；对于 $h$ 在 $\frakg$ 上定义的 Hodge 结构，这两个切空间分别具有 Hodge 型 $(-1,1)$ 与 $(1,-1)$。

我们有时写 $\frakp^+=\frakp^+_h$、$\mathfrak{k}=\mathfrak{k}_h$，以强调对 $h$ 的选取。

对称空间 $X$ 有 $G(\R)$-不变度量，除相差一个正实数倍外唯一，且它下降为 $S(G,X)$ 上的度量。若 $G/Z$ 各向异性，则 $S(G,X)$ 紧，其 de Rham 上同调可用关于该不变度量的调和形式来计算。Matsushima 公式把这一计算重新解释为自守形式空间 $\mathcal A(G)$ 的相对李代数上同调。

设 $\rho\colon G\to\mathrm{Aut}(V)$ 为有限维线性表示，定义在子域 $E\subset\C$ 上。设 $\tilde V$ 为 $S(G,X)$ 上相应的、系数在 $E$ 中的局部系统：
\[
\tilde V=\varprojlim_{K_f}G(\Q)\backslash X\times G(\mathbb{A}_f)\times V/K_f.
\]
为使 $\tilde V$ 成为局部系统数系统，需要假设：对所有充分小的 $K_f$，$Z(\Q)\cap K_f$ 在 $Z$ 中的 Zariski 闭包含于 $\rho$ 的核中；以下将不再说明地作此假设。记 $V_\C=V\otimes_E\C$，并同法定义 $\tilde V_\C$。

所谓 $G(\mathbb{A})$ 的容许不可约表示，我们一律指（在取定 $K_\infty$ 之下）不可约的 $(\frakg,K_\infty)\times G(\mathbb{A}_f)$-表示。

\begin{thm}{7.2.1}（Matsushima 公式）\label{thm:a6-04-13}
设 $G/Z$ 各向异性。则存在 $G(\mathbb{A}_f)$-表示的典范同构
\[
H^*(S(G,X),\tilde V_\C)\xrightarrow{\sim}\bigoplus_\pi m(\pi)H^*(\frakg,\mathfrak{k};\pi_\infty\otimes V_\C)\otimes\pi_f.
\]
\end{thm}

这里 $\pi$ 取遍 $G(\mathbb{A})$ 的容许不可约表示，且有可数直和分解
\begin{equation}\tag{7.2.3}\label{eq:a6-04-13}
\mathcal A(G)\xrightarrow{\sim}\bigoplus_\pi m(\pi)\pi
\end{equation}
其中 $m(\pi)$ 是一个（整）重数，而 $\pi\xrightarrow{\sim}\pi_\infty\otimes\pi_f$ 是分解，$\pi_\infty$（相应地 $\pi_f$）是不可约 $(\frakg,K_\infty)$-模（相应地 $G(\mathbb{A}_f)$-表示）。记号 $H^*(\frakg,\mathfrak{k};-)$ 表示相对李代数上同调。

这一结果的标准参考是 Borel 与 Wallach 的书 [BW00]，定义与证明均可参看该书。Matsushima 公式的意义在于：它把上同调的计算分成一个整体部分（对应重数 $m(\pi)$）和一个只依赖于 $\pi_\infty$ 的局部部分。我们只考虑使 $\pi_\infty$ 缓增的那些 $\pi$；此时上同调 $H^*(\frakg,\mathfrak{k};\pi_\infty\otimes V_\C)$ 在 [BW00] 中已完全算出，我们在下一节照抄答案。

当 $G/Z$ 非各向异性时，结果要复杂得多，证明也困难得多。特别地，不再有 (7.2.3) 那样的直和分解。尽管如此，Franke [Fra98]（基于 Borel 早期的工作）证明了定理 7.2.1 的如下类比：
\begin{equation}\tag{7.2.4}\label{eq:a6-04-14}
H^*(S(G,X),\tilde V_\C)\xrightarrow{\sim}H^*(\frakg,\mathfrak{k};\mathcal A(G)\otimes V_\C).
\end{equation}

\ysrc{517}

\subsubsection{尖点上同调}

我们定义\emph{尖点上同调}
\[
H^*_0(S(G,X),\tilde V_\C)\subset H^*(S(G,X),\tilde V_\C)
\]
为 $H^*(\frakg,\mathfrak{k};\mathcal A_0(G)\otimes V_\C)$ 在同构 (7.2.4) 之下于 $H^*(\frakg,\mathfrak{k};\mathcal A(G)\otimes V_\C)$ 中的像。Borel 证明：从尖点形式的相对李代数上同调到上同调的映射是单射。于是，若记
\begin{equation}\tag{7.2.5}\label{eq:a6-04-15}
\mathcal A_0(G)\xrightarrow{\sim}\bigoplus_\pi m_0(\pi)\pi,
\end{equation}
其中 $m_0(\pi)$ 表示 $\pi$ 在尖点谱中的重数，我们就得到 Matsushima 公式的类比：
\begin{equation}\tag{7.2.6}\label{eq:a6-04-16}
H^*_0(S(G,X),\tilde V_\C)\xrightarrow{\sim}\bigoplus_\pi m_0(\pi)H^*(\frakg,\mathfrak{k};\pi_\infty\otimes V_\C)\otimes\pi_f.
\end{equation}

\begin{rem}{7.2.2}\label{thm:a6-04-14}
在 (7.2.6) 与定理 7.2.1 中，相对李代数上同调其实应换成 $(\frakg,K^c_\infty)$-上同调，其中 $K^c_\infty$ 现为极大紧子群。这里所陈述的版本只在 $K^c_\infty$ 连通且 $G$ 的中心有限时才正确。本文的讨论可以轻易修改，以处理 $K^c_\infty$ 不连通或 $G$ 的中心所带来的麻烦。
\end{rem}

我们固定 $G(\mathbb{A}_f)$ 的一个表示 $\pi_f$，并令 $[\pi_f]_\infty$ 记满足 $\pi=\pi_\infty\otimes\pi_f\subset\mathcal A_0(G)$ 的 $G(\R)$ 的不可约表示之集。我们作如下假设。

\begin{hyp}{7.2.3}\label{thm:a6-04-15}
每个 $\pi_\infty\in[\pi_f]_\infty$ 都是缓增的。
\end{hyp}

我们可以定义
\[
H^*_0(S(G,X),\tilde V_\C)[\pi_f]=\mathrm{Hom}_{G(\mathbb{A}_f)}\bigl(\pi_f,H^*_0(S(G,X),\tilde V_\C)\bigr).
\]
由 (7.2.6)，我们便有
\begin{equation}\tag{7.2.7}\label{eq:a6-04-17}
H^*_0(S(G,X),\tilde V_\C)[\pi_f]\xrightarrow{\sim}
\bigoplus_{\pi_\infty\in[\pi_f]_\infty}m_0(\pi)H^*(\frakg,\mathfrak{k};\pi_\infty\otimes V_\C).
\end{equation}

\subsubsection{缓增情形的相对李代数上同调}

给定 $V_\C$ 时，使相对李代数上同调非平凡的缓增 $\pi_\infty$ 之集，在 [BW00, §3] 中已确定，尤其见其中的定理 3.3 与定理 5.1。每个这样的 $\pi_\infty$ 都同构于形如 $I_{P,\sigma,0}$ 的抛物诱导表示（采用 [BW00] 的记号），其中 $P$ 是 $G(\R)$ 的基本抛物子群，其 Langlands 分解为 ${}^0MAN$；$\sigma$ 是 ${}^0M$ 的离散系列表示，其无穷小特征由 $V$ 的最高权确定：
\[
\chi_\sigma=-\chi_{-s(\rho+\lambda(V))|_{\mathfrak b}}
\]
（见 [BW00, Theorem 3.3 (2)]），其中 $\lambda(V)$ 是 $V$ 的最高权，$\mathfrak b$ 是 $\mathrm{Lie}({}^0M)$ 的紧 Cartan 子代数，$s$ 是 Weyl 群中一个（唯一确定的）长度为 $\ell(s)=\dim N/2$ 的元素。此外，相对李代数上同调 $H^*(\frakg,\mathfrak{k};\pi_\infty\otimes V_\C)$ 在所引的两条定理中已显式确定。我们令
\[
q(G)=\frac{\dim(G)-\dim(K^c_\infty)}{2},\qquad
\ell_0(G)=\mathrm{rank}(G)-\mathrm{rank}(K^c_\infty),
\]
\[
q_0(G)=q(G)-\frac{\ell_0(G)}{2}.
\]
不变量 $\ell_0$ 等于 $\dim A$，其中 $A$ 是 $G(\R)$ 中极大分裂环面的分裂部分。于是 $q_0(G)$ 是整数（参看 [BW00, §4.3]），并且有
\begin{equation}\tag{7.2.8}\label{eq:a6-04-18}
H^{q_0(G)+j}(\frakg,K^c_\infty;\pi_\infty\otimes V_\C)\xrightarrow{\sim}
H^{q({}^0M)}\bigl({}^0\frakm,K_P;\sigma\otimes W_{s(\lambda+\rho)-\rho}\bigr)\otimes\wedge^j(\mathfrak{a}^*).
\end{equation}

\ysrc{518}

这里，$\mathfrak{a}^*$ 是分裂分量 $A$ 的李代数 $\mathfrak{a}$ 的线性对偶，$W_{s(\lambda+\rho)-\rho}$ 是 ${}^0M$ 的具有所指最高权的不可约有限维表示，而 $K_P$ 是极大紧子群 $K^c_\infty\cap P(\R)=K^c_\infty\cap{}^0M(\R)$。此外，${}^0M$ 是一个有离散系列的群，$q({}^0M)$ 是其相伴对称空间的维数。

特别地，由于 $\ell_0(G)=\dim A$，我们有
\begin{equation}\tag{7.2.9}\label{eq:a6-04-19}
H^q(\frakg,K^c_\infty;\pi_\infty\otimes V_\C)\ne0
\iff q\in[q_0(G),q_0(G)+\ell_0(G)].
\end{equation}

\subsubsection{外代数作用}

关于离散系列，Schmid 的一个基本事实是：(7.2.8) 中出现的空间
$H^{q({}^0M)}({}^0\frakm,K_P;\sigma\otimes W_{s(\lambda+\rho)-\rho})$ 是一维的。由此可得

\begin{cor}{7.2.4}\label{thm:a6-04-16}
$H^{q_0(G)+*}(\frakg,K^c_\infty;\pi_\infty\otimes V_\C)$ 自然地是外代数 $\wedge^\bullet(\mathfrak{a}^*)$ 上的自由秩一微分分次模。
\end{cor}

把这一结果与 (7.2.7) 结合，便得到如下整体推论，它是 Venkatesh 母题猜想的起点。

\begin{cor}{7.2.5}\label{thm:a6-04-17}
设 $\pi_f$ 具有如下性质：每个 $\pi_\infty\in[\pi_f]_\infty$ 都缓增。则
\[
H^{q_0(G)+\bullet}_0(S(G,X),\tilde V_\C)[\pi_f]
\]
是 $\wedge^\bullet(\mathfrak{a}^*)$ 上的自由微分分次模，其秩为
\[
m_0(\pi_f)=\sum_{\pi_\infty\in[\pi_f]_\infty}m_0(\pi).
\]
\end{cor}

Venkatesh 母题猜想的起点是如下观察：$\mathfrak{a}^*$ 的维数 $\ell_0$ 与一个假想地附着于 $\pi_f$ 的母题上同调群的秩一致。这一点在 [PV21, Ven19] 中有清楚的解释，我们在此复述。对于一个由 $G$ 的上同调尖点自守表示 $\pi$ 与 Langlands $L$-群 ${}^LG$ 的有限维表示 $\tau$ 构成的配对，Langlands 对应假想地给出一族母题 $M(\tau)_\pi$：其构造是把 $\tau$ 与 $\pi$ 的参数复合，而该参数被看作从 $\Q$ 的绝对伽罗瓦群的一个母题版本到代数群 ${}^LG$ 的同态。就本节及随后几节而言，我们只需要对偶伴随表示 $\tau=\mathrm{Ad}^*$ 的情形；此时记 $M_\pi=M(\mathrm{Ad})_\pi$。一般地，这些母题取值于某个数域 $E_\pi$，我们暂且忽略这一点。

附着于该假想母题的是母题上同调群 $H^1_{\mathrm{mot}}(\Q,M_\pi(1))$。它应当是某个 $K$-理论群张量 $\Q$ 之后的直和项，并且猜想为有限维 $\Q$-向量空间（更一般地，$E_\pi$-向量空间）；我们把这个空间记作 $\varpi_\pi$（读作「Va」，以纪念 Venkatesh）。虽然 $M_\pi$ 及其母题上同调都不可及，但若 $M_\pi$ 是 $\Q$ 上某个光滑射影代数簇的上同调中被对应截出的一个实际片段，那么存在调节子映射，比较 $\varpi_\pi\otimes C$ 与若干上同调实现空间，其中 $C$ 取遍可能的系数域
